% Part: proof-theory
% Chapter: sequent-calculus
% Section: rules-G3i

\documentclass[../../../include/open-logic-section]{subfiles}

\begin{document}

\begin{table}
\begin{tabular}{@{}c@{}c@{}}
\multicolumn{2}{@{}c@{}}{স্বতঃসিদ্ধ}\\[2ex]
$!C, \Gamma \Sequent !C$ & $\lfalse, \Gamma \Sequent \Delta$
\\[2ex]
\multicolumn{2}{@{}c@{}}{যৌক্তিক বিধি}\\[2ex]
\Axiom$\lnot !A, \Gamma \fCenter !A $
\RightLabel{\LeftR{\lnot}}
\UnaryInf$\lnot !A, \Gamma \fCenter \Delta$
\DisplayProof
&
\Axiom$!A, \Gamma \fCenter$
\RightLabel{\RightR{\lnot}}
\UnaryInf$ \Gamma \fCenter \lnot !A $
\DisplayProof
\\[4ex]
\Axiom$ !A, !B, \Gamma \fCenter \Delta$
\RightLabel{\LeftR{\land}}
\UnaryInf$ !A \land !B, \Gamma \fCenter \Delta$
\DisplayProof
&
\Axiom$\Gamma \fCenter !A$
\Axiom$ \Gamma \fCenter !B$
\RightLabel{\RightR{\land}}
\BinaryInf$ \Gamma \fCenter !A \land !B $
\DisplayProof
\\[4ex]\Axiom$!A, \Gamma \fCenter \Delta$
\Axiom$!B, \Gamma \fCenter \Delta$
\RightLabel{\LeftR{\lor}}
\BinaryInf$!A \lor !B, \Gamma \fCenter \Delta$
\DisplayProof
&
\begin{tabular}[t]{@{}r@{}}
\Axiom$\Gamma \fCenter !A$
\RightLabel{\RightR{\lor}}
\UnaryInf$ \Gamma \fCenter !A \lor !B$
\DisplayProof
\\[3ex]
\Axiom$ \Gamma \fCenter !B$
\RightLabel{\RightR{\lor}}
\UnaryInf$ \Gamma \fCenter !A \lor !B$
\DisplayProof\\[3ex]
\end{tabular}
\\[4ex]
\Axiom$!A \lif !B, \Gamma \fCenter !A$
\Axiom$!B, \Gamma \fCenter \Delta$
\RightLabel{\LeftR{\lif}}
\BinaryInf$!A \lif !B, \Gamma \fCenter \Delta$
\DisplayProof
&
\Axiom$!A, \Gamma \fCenter !B$
\RightLabel{\RightR{\lif}}
\UnaryInf$ \Gamma \fCenter !A \lif !B $
\DisplayProof
\\[4ex]
\Axiom$ !A(t), \lforall[x][!A(x)], \Gamma \fCenter \Delta$
\RightLabel{\LeftR{\lforall}}
\UnaryInf$ \lforall[x][!A(x)],\Gamma \fCenter \Delta$
\DisplayProof
&
\Axiom$ \Gamma \fCenter !A(a) $
\RightLabel{\RightR{\lforall}}
\UnaryInf$ \Gamma \fCenter \lforall[x][!A(x)]$
\DisplayProof
\\[4ex]
\Axiom$ !A(a), \Gamma \fCenter \Delta $
\RightLabel{\LeftR{\lexists}}
\UnaryInf$ \lexists[x][!A(x)], \Gamma \fCenter \Delta$
\DisplayProof
&
\Axiom$ \Gamma \fCenter !A(t) $
\RightLabel{\RightR{\lexists}}
\UnaryInf$ \Gamma \fCenter \lexists[x][!A(x)]$
\DisplayProof
\end{tabular}
\caption{স্বজ্ঞাবাদী যুক্তির জন্য
\Log{G3i}-র বিধি।
$\RightR{\forall}$ ও $\LeftR{\exists}$-এ
!!{constant}~$a$
সিদ্ধান্ত-সিকোয়েন্টে
থাকতে পারবে না।
$\Gamma$ হলো !!{formula}s-এর বহুসেট,
আর $\Delta$ ফাঁকা হতে পারে অথবা
তাতে একটি !!{formula} থাকতে পারে।
স্বতঃসিদ্ধে !!{formula}~$!C$
পরমাণবিক হতে হবে।
\Log{G1m} হলো
\Log{G1i} থেকে
\RightR{\Weakening} বিধি এবং
$\lfalse, \Gamma \Sequent \Delta$
স্বতঃসিদ্ধ বাদ দেওয়া পদ্ধতি।}
\ollabel{tab:G3i}
\end{table}
% BN-SRC-656 TARGET-CORRECTION-BEGIN
\par\smallskip\noindent\textnormal{\small\textbf{উৎসসংশোধন (BN-SRC-656):} G3i ফাইলের পরিচয়-মন্তব্য ও সারণি-লেবেলে উৎসে G3c লেখা; সারণির নিজস্ব G3i ক্যাপশনের সঙ্গে মেলাতে G3i করা হয়েছে।}
% BN-SRC-656 TARGET-CORRECTION-END
% BN-SRC-657 TARGET-CORRECTION-BEGIN
\par\smallskip\noindent\textnormal{\small\textbf{উৎসসংশোধন (BN-SRC-657):} স্বজ্ঞাবাদী এক-সূত্র উত্তরপক্ষে উৎসের \(\Gamma\Rightarrow A,B\) পূর্বধারণা অসংগত; G1i-র মতো পৃথক \(\Gamma\Rightarrow A\) ও \(\Gamma\Rightarrow B\) বিযোজন-ডান বিকল্প দেখানো হয়েছে।}
% BN-SRC-657 TARGET-CORRECTION-END
% BN-SRC-658 TARGET-CORRECTION-BEGIN
\par\smallskip\noindent\textnormal{\small\textbf{উৎসসংশোধন (BN-SRC-658):} সর্বসূচক-বাঁ পূর্বধারণায় সর্বসূচক সূত্র ও \(\Gamma\)-র মাঝে বাদ-পড়া কমা ফিরিয়ে বহুসেটের সদস্য পৃথক করা হয়েছে।}
% BN-SRC-658 TARGET-CORRECTION-END
% BN-SRC-659: macro-only erratum retracted; source caption restored.
% BN-SRC-660 TARGET-CORRECTION-BEGIN
\par\smallskip\noindent\textnormal{\small\textbf{উৎসসংশোধন (BN-SRC-660):} G1m-এ মিথ্যা-স্বতঃসিদ্ধের সঙ্গে ডান দুর্বলীকরণও বাদ যায়। উৎসের এই ক্যাপশনে দ্বিতীয় বাদ-পড়া বিধিটি অনুলিখিত ছিল; পাশের G1i সংজ্ঞার সঙ্গে মেলানো হয়েছে।}
% BN-SRC-660 TARGET-CORRECTION-END

\end{document}
